\documentclass[book.tex]{subfiles}
\begin{document}

\chapter{Deep Tensor Neural Networks}
\label{chap:dtnn}
\noindent\rule{11cm}{0.4pt} \\
\textbf{Prerequisites:} \autoref{chap:graphconvs}, \\
\textbf{Difficulty Level:} ** \\
\noindent\rule{11cm}{0.4pt}
\vspace{5mm}

Deep Tensor Neural Networks (DTNN) are adaptable extensions of the Coulomb matrix featurizer.\cite{schutt2017quantum} The core idea is to directly use nuclear charge (atom number) and the distance matrix to predict energetic, electronic or thermodynamic properties of small molecules. To build a learnable system, the model first maps atom numbers to trainable embeddings (randomly initialized) as atomic features. Then each atomic feature $a_i$ is updated
based on distance info $d_{ij}$ and other atomic features $a_j$.
Comparing with Weave models, DTNNs share the same idea in terms of updating based on both atomic and pair features, while the difference is using physical distance instead of graph distance. Note that the use of 3D coordinates to calculate physical distances limits DTNNs to quantum mechanical (or perhaps biophysical) datasets. We reimplement the model proposed by Sch{\"u}tt et al.\cite{schutt2017quantum} in a more generalized fashion. Atom numbers and a distance matrix are calculated by RDKit,\cite{landrum2006rdkit} using the Coulomb matrix featurizer. After embedding atom numbers into feature vectors $a_i$, we update $a_i$ in each convolutional layer by adding the
outputs from all network layers which use $d_{ij}$ and $a_j$ ($i$ $\neq$ $j$) as
input. After several layers of convolutions, all atomic features are summed together to form molecular features, used for classification and regression tasks.


\end{document}